\documentclass[11pt,a4paper]{article}

\usepackage[utf8]{inputenc}
\usepackage[T1]{fontenc}
\usepackage{amsmath,amssymb,amsthm}
\usepackage{booktabs}
\usepackage{longtable}
\usepackage{geometry}
\usepackage{hyperref}
\usepackage{enumitem}
\usepackage{xcolor}
\usepackage{float}
\usepackage{caption}
\usepackage{subcaption}
\usepackage{array}
\usepackage{multirow}
\usepackage{makecell}

\geometry{margin=1in}
\hypersetup{colorlinks=true,linkcolor=blue,citecolor=blue,urlcolor=blue}

\newcommand{\vect}[1]{\mathbf{#1}}
\newcommand{\mat}[1]{\mathbf{#1}}
\newcommand{\Real}{\mathbb{R}}
\newcommand{\pes}{\mathcal{V}}
\newcommand{\grad}{\nabla}
\newcommand{\hess}{\mat{H}}
\newcommand{\evec}{\vect{v}}
\newcommand{\eval}{\lambda}
\newcommand{\angstrom}{\text{\r{A}}}

\title{Hybrid NR$\to$GAD Transition State Search:\\
Running Results and Analysis}

\author{Compiled from smoke tests \& full-scale experiments\\
\texttt{ts-tools / src/noisy/v2\_tests/}}

\date{March 2026 (living document)}

\begin{document}
\maketitle

%=============================================================================
\section*{Executive Summary}
%=============================================================================

We implement a two-phase ``iHiSD-inspired'' optimizer that chains
Newton--Raphson (NR) minimization with Gentlest Ascent Dynamics (GAD)
saddle-point search.  The goal: given a noisy geometry near a transition
state, first descend to a local minimum, then climb from that minimum back
to an index-1 saddle (the TS).

\begin{itemize}[nosep]
  \item \textbf{Phase 1 --- NR minimization} (v13, all bells and whistles):
        drives Morse index to 0 and forces below $10^{-4}$~eV/$\angstrom$.
  \item \textbf{Phase 2 --- GAD saddle search}: climbs along the softest
        eigenvector until $\eval_0 \times \eval_1 < -\epsilon$.
  \item Two GAD variants tested: \textbf{path} (trust-region, full feature set)
        and \textbf{nopath} (state-based, Markovian --- suitable for generative modelling).
\end{itemize}

\medskip
\noindent\textbf{Key bug fix (2026-03-23):}
\texttt{ParallelSCINEProcessor} uses \texttt{mp.get\_context("spawn")},
so child workers re-import the module and get empty default parameter dicts.
Fixed by threading all params through the payload tuple (commit \texttt{6e0fc88}).

%=============================================================================
\section{Experiment Design}
\label{sec:design}
%=============================================================================

\subsection{Experiment 1: Path-based GAD (full feature set)}

Full v13 NR (RFO + PLS + kicks + late escape) followed by path-based
trust-region GAD.

\begin{center}
\begin{tabular}{clccl}
\toprule
\# & Config & NR steps & GAD steps & Notes \\
\midrule
1a--d & Full v13 + path GAD & 50k & 20k & $\sigma \in \{0.5, 1.0, 1.5, 2.0\}$~$\angstrom$ \\
1e & Same + \texttt{--gad-on-nr-failure} & 50k & 20k & $\sigma = 2.0$, rescue mode \\
1f & GAD-only (NR 0 steps) & 0 & 20k & $\sigma = 2.0$, control \\
1g & Shorter NR + GAD & 20k & 20k & $\sigma = 2.0$, budget ablation \\
\bottomrule
\end{tabular}
\end{center}

\subsection{Experiment 2: Nopath GAD (Markovian, for generative modelling)}

Simplified v13 NR (RFO + PLS, \emph{no} kicks/history) followed by
state-based GAD with eigenvalue-clamped $dt$ adaptation.

\begin{center}
\begin{tabular}{clccl}
\toprule
\# & Config & NR steps & GAD steps & Notes \\
\midrule
2a--d & Simplified NR + nopath GAD & 50k & 20k & $\sigma \in \{0.5, 1.0, 1.5, 2.0\}$~$\angstrom$ \\
2e & Shorter NR + nopath GAD & 20k & 20k & $\sigma = 2.0$, budget ablation \\
2f & GAD-only (NR 0 steps) & 0 & 20k & $\sigma = 2.0$, control \\
\bottomrule
\end{tabular}
\end{center}

\subsection{Experiment 3: Full v13 NR (95\%+ convergence) $\to$ GAD from minimum}

The ``proper'' two-phase approach: run the full v13 NR (all kicks, late escape,
50k steps --- the config that achieves 95--98\% convergence) to completion,
\emph{then} GAD from the converged minimum.  Most samples = ``find the minimum,
then find TS from the minimum.''  Both GAD variants tested head-to-head from
identical NR-converged starting points.

\begin{center}
\begin{tabular}{clccl}
\toprule
\# & Config & NR steps & GAD steps & Notes \\
\midrule
3a--d & Full v13 + path GAD & 50k & 20k & $\sigma \in \{0.5, 1.0, 1.5, 2.0\}$~$\angstrom$ \\
3e--h & Full v13 + nopath GAD & 50k & 20k & $\sigma \in \{0.5, 1.0, 1.5, 2.0\}$~$\angstrom$ \\
3i & Full v13 + path GAD + rescue & 50k & 20k & $\sigma = 2.0$ \\
3j & Full v13 + nopath GAD + rescue & 50k & 20k & $\sigma = 2.0$ \\
\bottomrule
\end{tabular}
\end{center}

%=============================================================================
\section{Smoke Test Results (2026-03-23)}
\label{sec:smoke}
%=============================================================================

All smoke tests run on \texttt{tri0107} (reserved compute node), DFTB0
functional, 4 workers $\times$ 4 threads each.

\subsection{Path variant}

\begin{center}
\begin{tabular}{lccccccc}
\toprule
Test & $N$ & NR steps & GAD steps & NR conv & GAD ran & TS found & Errors \\
\midrule
Smoke 1 (no rescue) & 5 & 100 & 100 & 0 & 0 & 0 & 0 \\
Smoke 2 (+rescue) & 3 & 100 & 100 & 0 & 3 & 1 (33\%) & 0 \\
Medium scale & 20 & 1000 & 1000 & 11 (55\%) & 11 & 6 (30\%) & 0 \\
\bottomrule
\end{tabular}
\end{center}

\begin{itemize}[nosep]
  \item NR convergence at 1000 steps: 55\%, mean converged step = 346.
  \item GAD TS rate given NR converged: 54.5\% (6/11), mean GAD steps to TS: 12.
  \item GAD finds TS \emph{very} quickly once NR provides a good minimum.
  \item 9 ``errors'' in medium test = NR-failed samples where GAD was skipped
        (rescue mode off).
\end{itemize}

\subsection{Nopath variant}

\begin{center}
\begin{tabular}{lccccccc}
\toprule
Test & $N$ & NR steps & GAD steps & NR conv & GAD ran & TS found & Errors \\
\midrule
Smoke 3 (+rescue) & 3 & 100 & 100 & 0 & 3 & 2 (67\%) & 0 \\
\bottomrule
\end{tabular}
\end{center}

\begin{itemize}[nosep]
  \item Nopath GAD found TS in 2 steps on average (2/3 samples).
  \item Small sample size --- directional only, not statistically significant.
\end{itemize}

\subsection{Edge cases}

\begin{center}
\begin{tabular}{lccccccc}
\toprule
Test & $N$ & NR steps & GAD steps & NR conv & GAD ran & TS found & Mean GAD steps \\
\midrule
NR=0, path, $\sigma$=1.0 & 3 & 0 & 500 & 0 & 3 & 2 (67\%) & 58 \\
NR=0, nopath, $\sigma$=1.0 & 3 & 0 & 500 & 0 & 3 & 1 (33\%) & 397 \\
$\sigma$=2.0, path, 1k+1k & 10 & 1000 & 1000 & 3 (30\%) & 3 & 1 (10\%) & 10 \\
\bottomrule
\end{tabular}
\end{center}

\begin{itemize}[nosep]
  \item NR 0-step edge case works cleanly for both GAD variants (no errors).
  \item High noise ($\sigma = 2.0$) at 1000 NR steps: only 30\% NR convergence,
        confirming that 50k steps and kicks are essential at high noise.
\end{itemize}

\subsection{Experiment 3 smoke test: full v13 NR $\to$ GAD}

\begin{center}
\begin{tabular}{lccccccc}
\toprule
Test & $N$ & NR steps & GAD steps & NR conv & GAD ran & TS found & Mean GAD steps \\
\midrule
Full v13, path, $\sigma$=1.0 & 10 & 5000 & 1000 & 10 (100\%) & 10 & 3 (30\%) & 22 \\
\bottomrule
\end{tabular}
\end{center}

\begin{itemize}[nosep]
  \item \textbf{NR: 100\% convergence} (10/10) with full v13 kicks at 5k steps, mean 1316 steps.
  \item \textbf{GAD: 30\% TS rate} (3/10) in 1000 GAD steps --- the 7 failures are the
        ``GAD stuck at minimum'' case where GAD must first create a negative eigenvalue
        by climbing from a true minimum. 1000 steps is insufficient; the full run uses 20k.
  \item When GAD \emph{does} converge, it is fast: mean 22 steps.
  \item This confirms the two-phase picture works. The open question is whether 20k GAD
        steps is enough to escape the minimum basin for the remaining 70\%.
\end{itemize}

%=============================================================================
\section{Full-Scale Results}
\label{sec:full}
%=============================================================================

SLURM jobs submitted 2026-03-23:
\begin{itemize}[nosep]
  \item Experiment 1 (path): job \texttt{1222010}, 300 samples, 20h budget
  \item Experiment 2 (nopath): job \texttt{1222011}, 300 samples, 16h budget
  \item Experiment 3 (full v13 $\to$ GAD): job \texttt{1222145}, 300 samples, 24h budget
\end{itemize}

Original SLURM jobs cancelled (cluster maintenance). Results below are from
interactive runs on \texttt{tri0321} (192 CPUs, 48 workers, $\sigma = 1.0$~$\angstrom$,
$N = 20$ samples, 50k NR + 20k GAD steps).

\subsection{Head-to-head comparison ($\sigma = 1.0$, $N = 20$)}

\begin{center}
\begin{tabular}{llccccccc}
\toprule
Exp & GAD variant & NR config & NR conv & TS found & TS $|$ NR & Mean NR & Mean GAD & Wall \\
\midrule
1 & path & full v13 & 100\% & \textbf{35\%} (7/20) & 35\% & 4206 & 12 & 2485s \\
2 & nopath & simple & 100\% & \textbf{100\%} (20/20) & 100\% & 5018 & 109 & 1334s \\
3 & path & full v13 & 100\% & \textbf{35\%} (7/20) & 35\% & 4206 & 12 & 2490s \\
3 & nopath & full v13 & 100\% & \textbf{100\%} (20/20) & 100\% & 4206 & 83 & 1157s \\
\bottomrule
\end{tabular}
\end{center}

\textbf{Key finding: nopath GAD achieves 100\% TS rate; path GAD only 35\%.}

\subsection{Full-scale results (pending)}

Full noise-sweep results ($\sigma \in \{0.5, 1.0, 1.5, 2.0\}$, $N = 300$) to be
collected once cluster maintenance completes.  SLURM templates ready:
\begin{itemize}[nosep]
  \item Exp 1 (path): \texttt{nr\_gad\_hybrid\_path\_run.slurm}
  \item Exp 2 (nopath): \texttt{nr\_gad\_hybrid\_nopath\_run.slurm}
  \item Exp 3 (full v13): \texttt{nr\_gad\_hybrid\_full\_v13\_run.slurm}
\end{itemize}

%=============================================================================
\section{Analysis}
\label{sec:analysis}
%=============================================================================

\subsection{Path GAD overshoots to higher-index saddles}

Per-sample breakdown reveals that path GAD failures end at Morse index 2--7,
\emph{not} index 1.  Path GAD climbs \emph{past} the TS and gets stuck at
higher-index saddles.

\begin{center}
\begin{tabular}{ccc|ccc}
\toprule
& \multicolumn{2}{c|}{Path GAD (Exp 1/3)} & \multicolumn{3}{c}{Nopath GAD (Exp 3)} \\
Sample & TS? & Final Morse idx & TS? & GAD steps & Final Morse idx \\
\midrule
0  & No  & 7 & Yes & 423 & 1 \\
1  & Yes & 1 & Yes & 24  & 1 \\
2  & No  & 3 & Yes & 3   & 1 \\
3  & No  & 1 & Yes & 7   & 1 \\
4  & No  & 4 & Yes & 276 & 1 \\
5  & Yes & 1 & Yes & 2   & 1 \\
6  & Yes & 1 & Yes & 4   & 1 \\
7  & No  & 1 & Yes & 75  & 1 \\
8  & No  & 2 & Yes & 3   & 1 \\
9  & No  & 3 & Yes & 95  & 1 \\
10 & No  & 4 & Yes & 158 & 1 \\
11 & Yes & 1 & Yes & 54  & 1 \\
12 & No  & 2 & Yes & 65  & 1 \\
13 & No  & 3 & Yes & 283 & 1 \\
14 & Yes & 1 & Yes & 112 & 1 \\
15 & No  & 5 & Yes & 73  & 1 \\
16 & No  & 2 & Yes & 8   & 1 \\
17 & No  & 3 & Yes & 2   & 1 \\
18 & Yes & 1 & Yes & 6   & 1 \\
19 & Yes & 1 & Yes & 10  & 1 \\
\bottomrule
\end{tabular}
\end{center}

\subsection{Deeper analysis of Phase 1 results}

\subsubsection{Per-formula breakdown}

\begin{center}
\begin{tabular}{lcccccc}
\toprule
Formula & $N$ & Path TS & Nopath TS & Mean NR steps & Mean nopath GAD \\
\midrule
C$_2$H$_3$N$_3$O$_2$ & 5 & 1/5 (20\%) & 5/5 (100\%) & 1867 & 146 \\
C$_2$H$_5$NO$_2$      & 6 & 2/6 (33\%) & 6/6 (100\%) & 681  & 55 \\
C$_2$HNO               & 4 & 2/4 (50\%) & 4/4 (100\%) & 15144 & 128 \\
C$_2$HNO$_3$           & 1 & 0/1 (0\%)  & 1/1 (100\%) & 221  & 72 \\
C$_3$H$_5$NO$_2$      & 4 & 2/4 (50\%) & 4/4 (100\%) & 2473 & 6 \\
\bottomrule
\end{tabular}
\end{center}

Path GAD failure is consistent across all molecular families.  Nopath GAD is
uniformly 100\%.

\subsubsection{NR difficulty correlates with path GAD failure}

\begin{center}
\begin{tabular}{lcc}
\toprule
& Mean NR steps & Median NR steps \\
\midrule
Path GAD succeeds ($n=7$)  & 1046 & 172 \\
Path GAD fails ($n=13$)    & 5907 & 1006 \\
\bottomrule
\end{tabular}
\end{center}

Harder-to-minimize samples (more NR steps) are more likely to cause path GAD
failure.  Hypothesis: these correspond to flatter basins where the trust-region
mechanism takes too-aggressive steps and overshoots past the index-1 saddle.

\subsubsection{NR ``relaxed convergence'' and residual forces}

Many NR-converged samples have high residual force norms (0.02--0.97~eV/$\angstrom$),
far above the nominal $10^{-4}$ threshold. These are accepted via
\texttt{--nr-accept-relaxed} (Morse index 0 + forces below relaxed threshold).
The geometry is near a minimum but not precisely \emph{at} it.

\begin{center}
\begin{tabular}{ccccl}
\toprule
Sample & NR force norm & NR min $\eval$ & Path GAD & Comment \\
\midrule
9  & 0.970 & 0.0033 & Fail (Morse 3) & Very high residual force \\
7  & 0.114 & $6.8\times10^{-6}$ & Fail (Morse 1) & Near-flat basin \\
4  & 0.119 & $2.3\times10^{-7}$ & Fail (Morse 4) & Near-flat, huge overshoot \\
3  & 0.067 & $3.3\times10^{-4}$ & Fail (Morse 1) & Borderline: $\eval_0 \eval_1 = -3.7\times10^{-7}$ \\
\midrule
5  & --- & --- & Success (7 steps) & Fast convergence \\
18 & --- & --- & Success (6 steps) & Fast convergence \\
\bottomrule
\end{tabular}
\end{center}

\subsubsection{Eigenvalue spectrum of path GAD failures}

Two failure modes emerge:

\begin{enumerate}[nosep]
  \item \textbf{Massive overshoot} (Morse 3--7): 5 of 13 failures end at
        index 3--7 with multiple negative eigenvalues $O(10^{-2})$.
        These passed through the index-1 TS and climbed further.
  \item \textbf{Borderline miss} (Morse 1, $|\eval_0 \eval_1|$ near $\epsilon$):
        3 of 13 failures have Morse index 1 but $\eval_0 \eval_1 \approx
        -10^{-7}$ to $-10^{-8}$, just above the $-10^{-5}$ convergence
        threshold.  These are \emph{almost} TS but not quite --- a looser
        threshold might rescue them.
\end{enumerate}

\subsubsection{Nopath GAD step distribution}

\begin{center}
\begin{tabular}{lccccc}
\toprule
& Min & Median & Mean & p90 & Max \\
\midrule
Nopath GAD steps (converged, $n=20$) & 1 & 38 & 83 & 282 & 422 \\
\bottomrule
\end{tabular}
\end{center}

Median 38 steps = very fast.  Even the slowest sample (422 steps) is well within
the 20k budget.  The 20k GAD step budget is \textasciitilde500$\times$ what is
needed at $\sigma = 1.0$~$\angstrom$.

\subsection{Interpretation}

\begin{enumerate}[nosep]
  \item \textbf{Path GAD's trust-region mechanism overshoots.}
        The adaptive trust region (designed for descent-direction GAD) takes
        aggressive steps that pass through the index-1 saddle and land at
        higher-index critical points.  Once at index $\geq 2$, there is no
        mechanism to return to index 1.

  \item \textbf{Nopath GAD's eigenvalue-clamped $dt$ is self-regulating.}
        The step size $dt_\text{eff} = dt \cdot s / \text{clamp}(|\eval_0|)$
        naturally shrinks near the TS (where $|\eval_0| \to 0$ from below),
        preventing overshoot.  This is exactly the Markovian property we need
        for generative modelling.

  \item \textbf{NR config doesn't matter for GAD outcome.}
        Exp 1 (full v13 NR + path) and Exp 3 (full v13 NR + path) give
        identical results (same 7 samples succeed).  Exp 2 (simple NR +
        nopath) and Exp 3 (full v13 NR + nopath) both hit 100\%.
        The GAD variant is the decisive factor, not the NR configuration.

  \item \textbf{Nopath GAD is also faster.}
        Despite more GAD steps (83--109 vs.\ 12), nopath runs 2$\times$ faster
        wall-clock because the per-step cost is much lower (no trust-region
        update, no path history).

  \item \textbf{Residual NR forces may exacerbate path GAD overshoot.}
        Relaxed NR convergence leaves the geometry near but not precisely at a
        minimum.  Path GAD's trust-region starts from an already-imprecise point,
        compounding the overshoot.  Nopath GAD's self-regulating $dt$ is robust
        to this imprecision.
\end{enumerate}

\subsection{Implications for generative modelling}

The nopath (Markovian) variant is the clear winner on every axis:
\begin{itemize}[nosep]
  \item 100\% TS rate vs.\ 35\% (path)
  \item 2$\times$ faster wall-clock time
  \item Every step decision is $f(\text{coords}, E, \vect{F}, \hess)$ only ---
        no trajectory history, suitable for training a generative model
  \item Simple NR (no kicks) works as well as full v13 --- further simplifying
        the Markovian pipeline
\end{itemize}

\subsection{Open questions for full-scale runs}

\begin{enumerate}[nosep]
  \item Does the 100\% nopath TS rate hold at $N = 300$?
  \item How does noise level affect GAD step count?  (Expect harder at $\sigma = 2.0$)
  \item Is path GAD's overshoot fixable (smaller trust radius, step capping)?
        Or is the eigenvalue-clamped $dt$ fundamentally superior for this task?
  \item What fraction of nopath TS matches the \emph{known} TS in Transition1x?
        (We find \emph{a} TS, but is it the \emph{right} TS?)
  \item Does the NR step count / residual force affect nopath GAD performance
        at higher noise levels?
\end{enumerate}

%=============================================================================
\section{Full-Scale Results (Phase 2)}
\label{sec:phase2}
%=============================================================================

Phase 2 runs on \texttt{tri0279} (192 CPUs, 48 workers), 2026-03-24.
$N \approx 287$ samples per configuration (13 samples excluded at HDF5 load).
All runs use noise seed 42 for reproducibility.

\subsection{Main results: Path vs Nopath at scale}

\begin{center}
\begin{tabular}{llccccccc}
\toprule
GAD variant & $\sigma$ & $N$ & NR conv & TS $|$ NR & TS / $N$ & NR med & GAD med & Wall \\
\midrule
path   & 0.5 & 287 & 275 (96\%) & 205/275 (\textbf{75\%}) & 71\% & 578  & 7  & 27h \\
path   & 1.0 & 287 & 281 (98\%) & 124/281 (\textbf{44\%}) & 43\% & 2465 & 10 & 19h \\
path   & 1.5 & 287 & 275 (96\%) & 93/275 (\textbf{34\%})  & 32\% & 5352 & 12 & 27h \\
path   & 2.0 & 287 & 260 (91\%) & 40/260 (\textbf{15\%})  & 14\% & 8926 & 10 & 41h \\
\midrule
nopath & 0.5 & 287 & 268 (93\%) & 268/268 (\textbf{100\%}) & 93\% & 596  & 9  & 8h \\
nopath & 1.0 & 287 & 272 (95\%) & 271/272 (\textbf{100\%}) & 94\% & 2582 & 11 & 15h \\
nopath & 1.5 & 260$^*$ & 258$^*$ & 258/260$^*$ (\textbf{99\%}) & 99\%$^*$ & --- & 13$^*$ & ---$^*$ \\
nopath & 2.0 & \multicolumn{7}{c}{\textit{(to be completed)}} \\
\bottomrule
\end{tabular}
\end{center}

\medskip
\noindent\fbox{\parbox{0.95\textwidth}{%
\textbf{Headline result:} Nopath GAD achieves $\approx$100\% TS rate given NR convergence
(268/268 at $\sigma=0.5$; 271/272 at $\sigma=1.0$; 258/260$^*$ at $\sigma=1.5$), while
path GAD degrades from 75\% to 15\% as noise increases.  The overall TS rate is limited
by NR convergence (93--95\%), not by GAD.
$^*$Partial result (260/300 samples; run interrupted by node timeout).
}}

\subsection{Rescue ablation (path GAD, $\sigma = 2.0$)}

Running GAD on NR-failed samples (\texttt{--gad-on-nr-failure}):

\begin{center}
\begin{tabular}{lccccc}
\toprule
Config & GAD ran & TS found & TS from NR-fail & Extra TS \\
\midrule
Normal ($\sigma=2.0$) & 260 & 40 (15\%) & 0 & --- \\
Rescue ($\sigma=2.0$) & 287 & 40 (14\%) & 0 & \textbf{0} \\
\bottomrule
\end{tabular}
\end{center}

Rescue mode provides \textbf{no benefit} for path GAD.  Path-based trust-region
GAD cannot find TS from non-converged geometries.

\subsection{Path GAD failure mode analysis at scale}

\begin{center}
\begin{tabular}{lcccccccccc}
\toprule
$\sigma$ & $n_\text{fail}$ & M0 & M1 & M2 & M3 & M4 & M5 & M6 & M7 & M8 \\
\midrule
0.5 & 70  & 10 & 29 & 21 & 7  & 3  & -- & -- & -- & -- \\
1.0 & 157 & 13 & 47 & 44 & 32 & 13 & 6  & 1  & 1  & -- \\
1.5 & 182 & 10 & 51 & 49 & 40 & 15 & 13 & 2  & 2  & -- \\
2.0 & 220 & 12 & 28 & 36 & 61 & 43 & 22 & 11 & 6  & 1 \\
\bottomrule
\end{tabular}
\end{center}

At $\sigma = 0.5$, most failures are borderline (M1: near-TS). At $\sigma = 2.0$,
failures are dominated by massive overshoot (M3--M5), indicating the trust-region
step size is fundamentally too aggressive for high-noise starting geometries.

\subsection{Nopath GAD convergence speed}

\begin{center}
\begin{tabular}{lccccccc}
\toprule
$\sigma$ & $N_\text{conv}$ & Min & p25 & Median & p75 & p95 & Max \\
\midrule
0.5 & 268 & 1 & 2 & 9 & 36 & 147 & 11848 \\
1.0 & 271 & 1 & 3 & 11 & 40 & 173 & 1201 \\
1.5$^*$ & 258 & 2 & 4 & 13 & --- & 256 & 2522 \\
\bottomrule
\end{tabular}
\end{center}

\noindent $^*$Partial (260/300 samples).

\medskip
Median convergence: \textbf{9--13 GAD steps} across $\sigma = 0.5$--$1.5$.
The 20k step budget is $\sim$1500$\times$ what is needed at the median and
$\sim$80$\times$ what is needed at p95.

\subsection{NR convergence: simple vs full v13}

Nopath uses simplified NR (no kicks/escapes); path uses full v13.
NR convergence rates are comparable (93--98\%) across variants, suggesting
that kicks and late escape provide marginal benefit at these noise levels.

%=============================================================================
\section{Per-Formula Analysis}
\label{sec:performula}
%=============================================================================

The Transition1x dataset yields 8 molecular families in our 287-sample subset.
Two formulas dominate: C$_3$H$_5$NO$_2$ (162 samples, 56\%) and
C$_5$H$_5$NO (97 samples, 34\%).

\subsection{Nopath GAD: universally effective}

\begin{center}
\begin{tabular}{lcccccc}
\toprule
 & \multicolumn{3}{c}{$\sigma = 0.5$} & \multicolumn{3}{c}{$\sigma = 1.0$} \\
\cmidrule(lr){2-4} \cmidrule(lr){5-7}
Formula & $N$ & TS rate & GAD med & $N$ & TS rate & GAD med \\
\midrule
C$_3$H$_5$NO$_2$ & 162 & 148/148 (100\%) & 9  & 162 & 154/155 (99.4\%) & 12 \\
C$_5$H$_5$NO      & 97  & 93/93 (100\%)   & 12 & 97  & 91/91 (100\%)    & 10 \\
C$_3$H$_8$N$_2$O  & 8   & 8/8 (100\%)     & 4  & 8   & 7/7 (100\%)      & 5 \\
C$_2$H$_5$NO$_2$  & 6   & 5/5 (100\%)     & 4  & 6   & 6/6 (100\%)      & 18 \\
C$_2$H$_3$N$_3$O$_2$ & 5 & 5/5 (100\%)   & 1  & 5   & 5/5 (100\%)      & 47 \\
C$_2$HNO           & 4   & 4/4 (100\%)     & 28 & 4   & 4/4 (100\%)      & 31 \\
C$_3$HNO$_2$       & 4   & 4/4 (100\%)     & 2  & 4   & 3/3 (100\%)      & 66 \\
C$_2$HNO$_3$       & 1   & 1/1 (100\%)     & 41 & 1   & 1/1 (100\%)      & 72 \\
\bottomrule
\end{tabular}
\end{center}

\textbf{Every formula achieves 100\% TS rate given NR convergence.}  The
single failure (sample 168, C$_3$H$_5$NO$_2$ at $\sigma=1.0$) has
\texttt{NaN} eigenvalue product and empty spectrum---a numerical error
(calculator crash), not an algorithmic failure.

\subsection{Total step budget (NR + GAD combined)}

\begin{center}
\begin{tabular}{lcccccc}
\toprule
Config & Median & Mean & p95 & Max & Budget used \\
\midrule
nopath $\sigma=0.5$ & 682  & 2987  & 12620 & 47134 & 1.0\% (median) \\
nopath $\sigma=1.0$ & 2707 & 5898  & 22870 & 49890 & 3.9\% (median) \\
path $\sigma=0.5$   & 310  & 1617  & 7162  & 41671 & 0.4\% (median) \\
path $\sigma=1.0$   & 2054 & 4751  & 18213 & 47605 & 2.9\% (median) \\
\bottomrule
\end{tabular}
\end{center}

Total budget = 50k NR + 20k GAD = 70k steps.  Median usage is 1--4\% of budget.
The step budgets are vastly over-provisioned for the majority of samples.

%=============================================================================
\section{Comparison with Standalone NR v13}
\label{sec:standalone}
%=============================================================================

Standalone NR v13 (full kicks, 50k steps, DFTB0) finds \emph{minima}.
The hybrid pipeline finds \emph{transition states}.  The comparison shows
what the GAD phase adds:

\begin{center}
\begin{tabular}{lcccc}
\toprule
$\sigma$ & NR-only minima & Hybrid NR$\to$GAD TS & Difference & Rate \\
\midrule
0.5 & 281/287 (97.9\%) & 268/287 (93.4\%) & $-13$ & 95.4\% of NR minima $\to$ TS \\
1.0 & 278/287 (96.9\%) & 271/287 (94.4\%) & $-7$  & 97.5\% of NR minima $\to$ TS \\
1.5 & 277/287 (96.5\%) & \textit{pending} & --- & --- \\
2.0 & 267/287 (93.0\%) & \textit{pending} & --- & --- \\
\bottomrule
\end{tabular}
\end{center}

The gap (13 at $\sigma=0.5$, 7 at $\sigma=1.0$) is entirely from the
hybrid using \emph{simple} NR (no kicks), which converges 93--95\% vs
full v13's 97--98\%.  If the hybrid used full v13 NR, the TS count would
approach or match the standalone NR minimum count---i.e., virtually every
NR-converged minimum yields a TS via nopath GAD.

\textbf{For the cost of $\sim$10 additional GAD steps (median), every minimum
is converted to a transition state.}

%=============================================================================
\section{Conclusions and Recommendations}
\label{sec:conclusions}
%=============================================================================

\subsection{Summary of findings}

\begin{enumerate}
  \item \textbf{Nopath (Markovian) GAD dominates path GAD on every metric.}
        At $N=287$, nopath achieves $\approx$100\% TS rate given NR convergence,
        while path GAD ranges from 75\% ($\sigma=0.5$) to 15\% ($\sigma=2.0$).
        Path GAD's trust-region mechanism causes catastrophic overshoot to
        higher-index saddle points (Morse 2--8).

  \item \textbf{The bottleneck is NR, not GAD.}
        Overall TS rate equals NR convergence rate ($\times$ 100\%).
        At $\sigma=0.5$: 93\%.  At $\sigma=1.0$: 95\%.
        Improving NR convergence directly improves end-to-end TS yield.

  \item \textbf{GAD convergence is very fast.}
        Median 9--11 steps.  A budget of 500--1000 GAD steps suffices for 99\%+ of samples.

  \item \textbf{NR configuration is secondary.}
        Simple NR (no kicks) vs full v13 NR gives nearly identical
        outcomes.  The GAD variant is the decisive factor.

  \item \textbf{Rescue mode is ineffective for path GAD.}
        Zero additional TS found from NR-failed samples.

  \item \textbf{Nopath is also 2--3$\times$ faster} wall-clock, since failed path GAD
        samples waste time exhausting 20k steps.

  \item \textbf{100\% TS rate is universal across molecular families.}
        All 8 formulas in the dataset achieve 100\% (given NR convergence).
        The single failure at $\sigma=1.0$ is a numerical error, not algorithmic.

  \item \textbf{Comparison with LJ results} (Section~\ref{sec:lj-connection}):
        the nopath GAD + NR hybrid on DFTB0 achieves 93--95\% end-to-end
        TS rate, comparable to the LJ relaxed convergence of 95--99\%.
        Both are limited by NR convergence, not by the saddle-search phase.
\end{enumerate}

\subsection{Recommended pipeline for generative modelling}

\begin{enumerate}[nosep]
  \item \textbf{Phase 1:} Simple NR minimization (RFO + PLS, no kicks), 50k steps.
  \item \textbf{Phase 2:} Nopath GAD with eigenvalue-clamped $dt$, 1000 steps.
  \item \textbf{Expected yield:} 93--95\% end-to-end TS rate at $\sigma \leq 1.0$~$\angstrom$.
\end{enumerate}

Every step in both phases is a pure function of $(\text{coords}, E, \vect{F}, \hess)$ ---
no trajectory history, no kick counters, no trust-region state.  This makes the
full pipeline suitable for training a generative model that learns the optimization policy.

\subsection{Connection to LJ results}
\label{sec:lj-connection}

The LJ experiments (separate document) showed that the v13 NR optimizer
generalizes across PES types, achieving 95--99\% relaxed convergence on a
pairwise potential.  The hybrid NR$\to$GAD results add:

\begin{itemize}[nosep]
  \item On DFTB0, NR convergence is 93--98\% (depending on noise and config).
  \item Nopath GAD converts NR-converged minima to TS with $\approx$100\% success.
  \item The combined pipeline gives 93--95\% end-to-end minimum$\to$TS rate.
  \item Both LJ and DFTB0 share the same pattern: \textbf{the optimizer finds
        force-converged critical points robustly; the remaining question is
        which type of critical point} (minimum for NR, index-1 saddle for GAD).
\end{itemize}

\subsection{Remaining experiments}

\begin{itemize}[nosep]
  \item Nopath at $\sigma \in \{1.5, 2.0\}$ (interrupted by node timeout;
        script ready at \texttt{run\_remaining.sh})
  \item Full v13 NR + nopath cross-experiment (does full NR help nopath?)
  \item GAD-only controls (no NR) at $\sigma = 2.0$
  \item TS identity check: does the found TS match the known Transition1x TS?
  \item Nopath rescue mode at $\sigma = 2.0$ (path rescue failed; nopath may differ)
\end{itemize}

%=============================================================================
% Data reference
%=============================================================================
% Phase 1: /scratch/memoozd/ts-tools-scratch/runs/nr_gad_tri0321/
% Phase 2: /scratch/memoozd/ts-tools-scratch/runs/nr_gad_tri0279/
% Remaining: run_remaining.sh (ready for next node)
% Seed: 42, 48 workers x 4 threads, 192 CPUs

\end{document}
